\section{Hybrid Proof Sketch}
\label{sec:hybrids}

We sketch a sequence from the simulator of Section~\ref{sec:simulator}
to the real execution.  The main order is to restore the encrypted
states first, then the garblings containing the decryption keys, then
the private deliveries, and finally public sampling.  Each cryptographic
step is a family of one-challenge reductions: we finish that type of
replacement across the execution before proceeding to the next type.
We assume throughout that all committees satisfy their honesty
thresholds; choosing parameters to ensure this is separate from the
hybrid argument.

Until the last step, retain $\SimSetup$ and $\SimRO$ as algorithms
using the memoized service $\SimSample$.  The hybrids change that
service and the simulated party behaviour.  The entire public-sampling
computation, including the adversary's local evaluations, remains part
of each reduction.  We are not assuming an additional interface for
programming universal samples.

For the analysis, run the ideal computation inside the hybrid experiment.
Its internal record gives the reduction the actual inputs, outputs, and
states $\sigma_s$.  These values are not disclosed to the adversary or
environment.  This is proof bookkeeping, not extra information supplied
to the implementable simulator.

\paragraph{Nonce freshness.}
First establish that the service has not already answered the eventual
description $d_s$ before its non-nonce fields are frozen.  Work in an
auxiliary experiment stopped before the first honest nonce-committee
key is disclosed.  Message privacy of the nonce-share encryption hides
the honest shares, while threshold privacy hides $Q_s$ from the corrupt
shares.  The needed encryption privacy follows from RNCE security;
its dummy-message explanation also accounts for the simulator's
nonce-only keys.  Thus a premature query must guess a hidden uniform
nonce.  Count every service query, including queries made internally by
$\SimSetup$ and $\SimRO$, not only evaluations declared by the adversary.
Apply this argument up to the first premature hit along the nonce chain.

This stopped argument need not explain the substituted nonce packets
after their keys are revealed.  In the main hybrids, the nonce packets
are unchanged.  Robust reconstruction fixes the nonce uniquely, and
the remaining description fields cannot change after its release.
Keep the early-query abort in the games; its negligible probability
then transfers through the indistinguishable hybrids below.

\begin{description}
  \item[$H_0$: the simulated execution.]
    Start with the stated simulator: topology-only simulated garblings,
    non-committing deliveries, and state ciphertexts
    $c_s\sample\Enc(\statepk_{s+1},0^{|\sigma_s|})$.
    Selected sharings are completed only when the simulated input
    encoding is available.

  \item[$H_1$: restore the plaintext states.]
    Replace the state ciphertexts, one transfer at a time, by
    \[
      c_s\sample\Enc(\statepk_{s+1},\sigma_s).
    \]
    All garblings and the honest label deliveries remain simulated.
    For a target transfer, use the IND-CPA challenger's public key as
    $\statepk_{s+1}$ and call
    \[
      \Challenge(0^{|\sigma_s|},\sigma_s).
    \]
    The reduction knows $\sigma_s$ from the internal ideal computation.
    It never needs the challenge secret key: the successor's
    $\SimGarble$ uses only topology, and its $\SimEncode$ is given the
    prescribed output rather than decrypting the incoming ciphertext.
    No real garbling containing that key is present.
    This is the use of IND-CPA security from
    Definition~\ref{def:ind-cpa}.

  \item[$H_2$: choose hidden constants early.]
    Rewrite the experiment to choose each private input permutation
    $\pi_w$ and output pad $r_i$ when its round is prepared, and later use
    \[
      j=b\oplus\pi_w,\qquad z_i=y_i\oplus r_i.
    \]
    This is the same distribution as choosing uniform $j$ or $z_i$
    first and explaining the hidden value afterward.  Before their
    exposure or use, the simulated deliveries and garblings reveal
    neither $\pi_w$ nor $r_i$.  Earlier output-recipient corruption
    simply exposes the already-chosen pad.

    Also choose the outgoing state-encryption coins $\rho_s$ at
    preparation rather than evaluation.  They have no earlier observable
    use while the garbling remains simulated.  This exact
    deferred-decisions change is made \emph{after} the IND-CPA sweep:
    that sweep therefore did not require knowledge of the challenger's
    encryption coins.  Each augmented circuit $C_s$ is now fully
    specified when its garbling is requested.

  \item[$H_3$: restore the garblings.]
    Replace the simulated garblings by real garblings, one round at a
    time, leaving all honest label-delivery ciphertexts simulated.
    For the target round $s$, generate its genuine state-transfer keys
    and construct $C_s$ with $\statesk_s$, its output pads, and
    $\rho_s$ hardcoded.  At sample preparation, call
    $\GarbleCircuit(C_s)$.  Once the complete input is fixed, call
    $\EncodeInput(\xi_s)$, where $\xi_s$ consists of the application
    bits, incoming state ciphertext, and successor public key.

    Since all state ciphertexts now carry the correct states,
    correctness of state decryption gives
    \[
      C_s(\xi_s)=
      \bigl(y_s\oplus r_s,\,
        \Enc(\statepk_{s+1},\sigma_s;\rho_s)\bigr),
    \]
    with $r_s$ the vector of output pads.  This is precisely the output
    previously supplied to $\SimEncode$.  The first round uses the
    prescribed empty state.

    The game returns only the selected input encoding.  That is enough:
    use $\CompleteShares$ to complete each selected committee's fixed
    corrupt shares to the returned label bundled with $\GarbleAux$,
    and explain the honest members' RNCE packets before their keys are
    posted.  Neither inactive labels nor the encoding key are needed
    by this reduction.
    This is the use of adaptive garbling privacy from
    Definition~\ref{def:adaptive-garbling}: inputs arrive after the
    garbling, and simulation reveals only topology rather than embedded
    gate constants.

    The garbling reduction knows $\statesk_s$ because it generates the
    state keys itself.  It is a different reduction from the earlier
    IND-CPA reduction; there is no requirement to retain that reduction's
    unknown challenge key.

  \item[$H_4$: prepare both committee sharings early.]
    All garblings are now real.  The reduction for the next step can
    generate them locally and retain their full encoding keys.
    For each wire, complete both committee sharings during sampling,
    using independent sharing randomness for the two labels and their
    bundled common material.  Condition on the corrupt shares that were
    fixed first.

    Perfect threshold privacy and
    Lemma~\ref{lem:sharing-completion} show that this has the same
    distribution as honest sharing, or as completing the selected branch
    later.  Unselected honest members remain silent.  This exact rewrite
    makes every delivery plaintext available at sample preparation.

  \item[$H_5$: restore genuine deliveries.]
    Apply receiver non-committing security, one delivery key at a time.
    For a target key, call $\GetPublicKey$ at registration,
    $\EncryptMessage$ with its actual packet when the designated sample
    is prepared, and $\CorruptReceiver$ when its key must be disclosed.
    For a selected committee member, this last call is before its
    posting; for an output recipient, it is at its chosen corruption.
    An inactive honest member need never be exposed.
    The returned receiver tape accounts for the full retained state,
    not merely a decryption key.

    If an output key is exposed before its designated sample, or is
    assigned only to a nonce committee, use $\EncryptMessage$ on a
    dummy message, discard that ciphertext without publishing it, and
    call $\CorruptReceiver$.  Later visible ciphertexts under this
    fixed key are ordinary encryptions.  Ciphertexts in other sampling
    descriptions are likewise generated locally.
    This is the use of Definition~\ref{def:receiver-nce}.
    It comes after the garbling sweep because the reduction now knows
    both labels and hence all share packets when $\EncryptMessage$
    must be called.

  \item[$H_6$: restore public universal sampling.]
    The service now generates an honest sample $d(r_d)$ with independent
    uniform randomness for each new description.  The early-query abort
    can be removed at negligible cost.  Apply
    Definition~\ref{def:adaptive-universal-sampler} to replace
    $\SimSetup$ and $\SimRO$ with real setup and the random oracle.
    Its indistinguishability clause hides this change, and its
    honest-sample consistency clause ensures that public evaluations
    agree with the samples used to run the parties.  Consistency can
    also be checked in the preceding hybrids and transferred by the
    same comparisons.

    At this point all sampled data, keys, and party messages have their
    real distributions.  Robust reconstruction and correctness of
    encryption and garbling identify the delivered outputs and state
    transitions with the real protocol.
\end{description}

Each sweep ranges over the polynomially many rounds or roles in the
particular efficient experiment being analysed.  This is not a
setup-time bound on the number of rounds.  The outline separates exact
distributional rewrites from cryptographic comparisons; formalizing the
one-challenge reductions and their losses is the remaining proof work.

\paragraph{Activation checks.}
The current input phase admits only the designated input providers,
each protected until its single irrevocable posting.  Under this
restriction there is no additional activation-forgery step in the
sequence above.  If other corrupt roles are allowed to submit
activations during that phase, one-wayness of the authorization images
must additionally exclude a valid preimage arriving before the honest
provider's post.  The first-valid-post rule handles subsequent exposure
in either case.
